% 08_biophysics.tex
% Ch.8: Protein backbone impedance

\chapter{Biophysics: Protein Folding}
\label{ch:biophysics}

\section{Protein Backbone: From Proton Radius to Folding}
\label{sec:protein_bridge}

The protein $S_{11}$ folding engine uses the C$\alpha$--C$\alpha$
virtual bond distance $d_0 \approx 3.80$~\AA\ as its fundamental
length scale.  This traces back to the AVE axioms via:
\begin{resultbox}{Proton-to-Backbone Length Scale}
\begin{equation}
    d_p = \frac{4\hbar}{m_p c} \approx 0.841\;\text{fm}
    \;\xrightarrow{\text{nuclear bond}}
    r_{\text{cov}}
    \;\xrightarrow{\text{backbone}}
    d_0 \approx 3.80\;\text{\AA}
\end{equation}
\end{resultbox}
The ratio $d_0 / a_0 \approx 7.18$ (where $a_0 = \ell_{\text{node}}
/\alpha$ is the Bohr radius) suggests a geometric origin in the
peptide repeat unit.  The amino acid impedance table $Z_{\text{topo}}$
assigns complex impedances (R + jX) to each of the 20 residues,
with $Q_{\text{backbone}} \approx 7$ suppressing reactive coupling.

\begin{figure}[H]
    \centering
    \includegraphics[width=0.95\textwidth]{protein_backbone_impedance.pdf}
    \caption{Protein backbone impedance from AVE axioms.
    \textbf{Left:} Full derivation chain from $\ell_{node}$ to
    the C$\alpha$--C$\alpha$ peptide bond length $d_0 = 3.80$~\AA.
    \textbf{Centre:} Amino acid impedance spectrum --- hydrophobic
    (gold), positive/inductive (blue), negative/capacitive (red).
    \textbf{Right:} Local inter-residue reflection coefficient
    $\Gamma$ along a sample peptide sequence; high $|\Gamma|$
    indicates folding-driving impedance mismatch.}
    \label{fig:protein_backbone}
\end{figure}

\section{Derivation Chain: From Lattice Pitch to Bond Length}

The C$\alpha$--C$\alpha$ virtual bond distance $d_0 \approx 3.80$~\AA\ is derived through a chain of AVE identities:

\begin{enumerate}
    \item \textbf{Lattice pitch:} $\ell_{node} = \hbar / (m_e c) \approx 3.862 \times 10^{-13}$ m.
    \item \textbf{Bohr radius:} $a_0 = \ell_{node} / \alpha \approx 5.292 \times 10^{-11}$ m (the characteristic extension of the unknot electron).
    \item \textbf{Covalent bond length:} C--N and C$\alpha$--C bonds are set by the overlap of atomic ground states at the equilibrium impedance-matching distance: $r_{cov} \approx 1.47$~\AA\ (C--N), 1.52~\AA\ (C--C), 1.33~\AA\ (C=O).
    \item \textbf{Backbone repeat:} The peptide unit (N--C$\alpha$--C--N) has a fixed geometry from the $sp^3$ tetrahedral bond angle ($\theta = 109.47^\circ$) and the planar amide constraint ($\omega = 180^\circ$). The resulting C$\alpha$--C$\alpha$ distance:
    \begin{equation}
        d_0 = \sqrt{r_{C\alpha C}^2 + r_{CN}^2 + 2\, r_{C\alpha C}\, r_{CN} \cos(180^\circ - \theta)} \approx 3.80 \text{ \AA}
    \end{equation}
\end{enumerate}

Every step uses constants from the AVE physics engine (Volume~I) or standard NERF rigid-body geometry. No free parameter is introduced at the biological scale.

\section{Amino Acid Impedance Classification}

The 20 canonical amino acids are classified by their sidechain impedance relative to the backbone baseline ($Z_{backbone} \approx 7$):

\begin{itemize}
    \item \textbf{Hydrophobic (low $Z$):} Ala, Val, Leu, Ile, Pro, Phe, Trp, Met --- compact sidechains with low reactive coupling. $|\Gamma| \approx 0$ (impedance-matched, minimal folding drive).
    \item \textbf{Polar/Charged (high $|Z|$):} Glu, Asp (capacitive), Lys, Arg (inductive) --- large $|\Gamma|$, strong folding drive at sequence positions where impedance mismatch is maximal.
    \item \textbf{Special:} Gly ($Z \approx 0$, maximum flexibility), Cys (disulphide bridge, reactive short-circuit between distant residues).
\end{itemize}

\section{Validation: Chignolin (CLN025)}

The 10-residue Chignolin mini-protein (CLN025: YYDPETGTWY) is the benchmark system for the AVE folding engine (Volume~V). Using zero adjustable parameters, the engine folds from an extended chain to a $\beta$-hairpin with:
\begin{resultbox}{Chignolin Validation}
\begin{equation}
    \text{RMSD}_{\text{backbone}} = 2.59 \text{ \AA} \quad \text{(sub-3 \AA: Level 6 Emergence)}
\end{equation}
\end{resultbox}

This validation uses the \textit{same three functions} as every other chapter in this volume: \texttt{impedance()}, \texttt{saturation\_factor()}, and \texttt{reflection\_coefficient()}. The protein backbone is not ``biology''---it is a transmission line.

